\section{Participation, Act, and the Limits of the Ancients}

When philosophy first woke in the Greek world, it did not begin with the question of why there is something rather than nothing. It began with a more immediate puzzle: how anything can remain intelligible across ceaseless change. Water freezes to ice and boils to vapor; wood burns to ash; bodies grow, sicken, and decay. What abides beneath the flux so that speech and knowledge are possible?

The earliest Ionian thinkers answered in terms of a single material substrate---water, air, fire, or an indefinite elemental mass---imagining that every observable change was merely the thickening or thinning, combining or separating, of an eternal bodily stuff. Everything was already there; the only problem was how its surface qualities shifted. That first stage of inquiry was indispensable, yet it remained trapped inside the horizon of accidental alteration. It took Plato and Aristotle to see that the visible world cannot even account for its own forms and motions without pointing beyond the changing realm altogether.

\subsection{Plato: Participation and the Good Beyond Being}

Plato broke decisively with material reductionism by observing that no sensible thing is identical with the perfection it displays. In the \emph{Phaedo}, Socrates points out that two sticks or stones may look roughly equal to one another, yet they can appear unequal from a different angle or be chipped away tomorrow; they fall short of Equality itself, which never alters and never admits its opposite. Similarly, a beautiful vase fades and breaks, a just city slides into tyranny, and a living body dies. Sensible realities possess beauty, justice, equality, or unity only in a partial, shifting, derivative manner.

Plato named this derivative possession \emph{participation} (\emph{methexis}). Whatever is beautiful without being Beauty itself, or just without being Justice itself, is what it is only by sharing in an unchanging Form that is that perfection by its very nature. The logic of participation already contains a structural asymmetry: the many changing instances depend on the one unchanging archetype for their intelligibility and character, while the archetype depends on none of its shifting copies.

In the \emph{Republic}, Plato pushed that hierarchy to its summit. Just as the sun in the visible realm not only makes objects visible to the eye but also supplies the light and warmth by which living things grow and flourish, so in the intelligible realm the Form of the Good stands at the apex of all Forms. It is not merely the source of truth and knowability for the mind; Plato explicitly describes the Good as the cause of the very being and essence of the intelligible Forms themselves, while standing beyond essence in dignity and power. Here Greek philosophy came close to a transcendent source of reality.

Yet when Plato turned in the \emph{Timaeus} to give an account of how the visible cosmos comes to be, the exact boundary of his metaphysics came into view. The Demiurge, or divine Craftsman, is good and free of all envy, desiring that all things should become as like himself as possible. Looking to the eternal intelligible model, he brings order out of disorder, persuading wild, wandering necessity to accept rational proportion. Even so, the Demiurge does not bring the underlying medium into existence. Alongside the eternal Forms and the good Craftsman stands a third, unoriginated principle: the \emph{chora} or Receptacle---an eternal, formless, spatial matrix and chaotic flux of pre-existing material necessity that resists rational form and limits what even divine wisdom can accomplish. The Platonic Craftsman shapes and orders; he does not bring matter itself out of nothing.

\subsection{Aristotle: Potency, Act, and the Unmoved Actuality}

Aristotle found Plato's separated Forms unable to explain the concrete reality of physical change: how can a static, separated blueprint sitting outside the material world actually move a seed to become a tree? In his \emph{Physics} and \emph{Metaphysics}, Aristotle forged the analytical tools that shaped Western philosophy for two millennia: the distinction between \emph{potency} (\emph{dynamis}) and \emph{act} (\emph{energeia}), and the architecture of fourfold causality.

For Aristotle, change is neither an illusion (as Parmenides held) nor sheer unintelligible flux (as Heraclitus suggested). Change is the actualization of a real potentiality in a real subject. Cold water is not actually hot, yet unlike a stone's inability to see, it possesses a real potency to become hot; fire, which is actually hot, reduces that potency to act. From this analysis in \emph{Metaphysics} IX, Aristotle established a foundational principle: actuality is prior to potentiality in definition, in knowledge, and in substance. Nothing in mere potency can actualize itself, because what is not yet actual lacks the very perfection it would need in order to bestow it. Every transition from potency to act therefore requires an already actual cause.

In \emph{Metaphysics} II (\emph{Alpha Elatton}), Aristotle then applied this causal logic to ordered series of causes. Whenever a sequence of causes stands in an essential order---a first cause, an intermediate cause, and a last effect---an infinite regress is impossible. Remove the first term, and there is no cause at all; every intermediate member transmits motion or causality only insofar as it receives it from what stands prior to it. Aristotle also noted in the same book a principle of metaphysical gradation that later scholasticism quoted repeatedly: whatever possesses a perfection in the highest degree is the cause of that perfection in all other things that share it, just as fire, being maximally hot, is the cause of heat in hot bodies, so that what is greatest in truth is also greatest in being.\endnote{Aquinas quotes this principle from Aristotle's \emph{Metaphysics} II directly in \emph{Summa Theologiae} I, q.\ 2, a.\ 3 (Fourth Way) and I, q.\ 44, a.\ 1: ``whatever is greatest in being and greatest in truth, is the cause of every being and of every truth; just as whatever is the greatest in heat is the cause of all heat.''}

In \emph{Metaphysics} XII (\emph{Lambda}), Aristotle carried the priority of act to its cosmic conclusion. Because physical motion and time, in Aristotle's view, are eternal and uninterrupted, there must exist an eternal, immaterial, wholly actual substance whose very essence is actuality (\emph{energeia}) without any admixture of potency or matter. Were this principle to contain even the slightest unrealized potency, it would remain capable of not acting, and cosmic motion could fail. Therefore the First Principle must be an Unmoved Mover, drawing the outermost celestial sphere into eternal circular motion not by mechanical pushing, but as an object of intellectual desire and final attraction---a living, eternal Mind whose sole activity is unbroken self-contemplation, thought thinking itself.

Aristotle thus showed by philosophical analysis that a world of potency and act cannot terminate in another moved mover or in an endless regress of dependent movers; it requires pure, immaterial actuality. Yet Aristotle's Unmoved Mover does not bring the cosmos into existence out of nothing. In his universe, prime matter, the celestial spheres, and the species of living things are co-eternal with the First Principle. The Unmoved Mover is the supreme final cause of cosmic motion and order, not the efficient cause of the existence of matter itself. Thinking only of itself, this eternal Mind neither originates primary matter nor freely wills particular things into existence.

\subsection{The Three-Stage Ascent to Universal Being}

More than fifteen centuries after Plato and Aristotle, St.\ Thomas Aquinas traced the entire arc of Greek metaphysics and diagnosed its exact ceiling in the \emph{Summa Theologiae}, examining how ancient philosophy advanced toward a universal cause of being.\endnote{St.\ Thomas Aquinas, \emph{Summa Theologiae} I, q.\ 44, a.\ 2. All quotations from the \emph{Summa Theologiae} follow the public-domain translation of the Fathers of the English Dominican Province.}

Aquinas describes the ancient philosophers advancing ``gradually, and as it were step by step,'' toward the truth. At the first stage, thinkers of ``grosser mind'' recognized no realities beyond sensible bodies. Supposing corporeal substance itself to be unoriginated, they sought causes only for \emph{accidental} changes---rarefaction and condensation, union and separation---assigning finite principles such as affinity, discord, or ordering intellect to rearrange what already existed.

At the second stage, philosophy achieved a decisive breakthrough. Thinkers grasped the distinction between substantial form and underlying matter, recognizing that bodies undergo not merely surface modifications but genuine substantial changes---when food becomes flesh, or elements combine into a new substance. To account for these substantial transmutations, they looked beyond particular colliding bodies to universal causes, ``such as the oblique circle, or ideas, according to Plato.'' Yet both Plato and Aristotle left primary matter itself untouched at this second stage: they ``imagined'' primary matter ``to be uncreated.'' Even when explaining why matter becomes \emph{this} kind of thing (human, oak, gold) rather than merely \emph{such} a colored or heated thing, they were still accounting for being only under a particular determination---``either as `this' or as `such'''---while treating the bare existence of the underlying world as an unexamined backdrop.

Then came the third and deepest ascent. Aquinas writes:

\begin{quote}
Then others there were who arose to the consideration of ``being,'' as being, and who assigned a cause to things, not as ``these,'' or as ``such,'' but as ``beings.'' Therefore whatever is the cause of things considered as beings, must be the cause of things, not only according as they are ``such'' by accidental forms, nor according as they are ``these'' by substantial forms, but also according to all that belongs to their being at all in any way. And thus it is necessary to say that also primary matter is created by the universal cause of things.\endnote{Aquinas, \emph{ST} I, q.\ 44, a.\ 2, corpus.}
\end{quote}

Between the second stage and the third, the horizon of inquiry shifts decisively. At the first stage, the question is why bronze is hot or cold (\emph{such}). At the second stage, the question is why matter has the form of bronze or of a living horse (\emph{this}). At the third stage, the question penetrates beneath every accidental quality and even beneath substantial species to ask why there is \emph{being at all}---why anything stands outside nothingness in any way whatever. Once that question is asked, nothing in the finite world can sit outside causal dependence: not form, not individual composite substances, and not even primary matter, which is pure potentiality and cannot exist apart from form yet falls wholly under the universal cause of being.

An objector in the Aristotelian tradition might reply that the ancients rightly accepted the axiom ``nothing is made from nothing.'' In \emph{ST} I, q.\ 45, a.\ 2, ad 1, Aquinas answers directly. The ancient philosophers were analyzing particular agents producing particular effects---a sculptor carving marble, a fire heating iron, nature generating an animal from parents. Every particular agent necessarily presupposes pre-existing matter on which to act, because a particular agent causes only a particular form or alteration, not universal being. Within the realm of physical becoming, ``nothing comes from nothing'' remains wholly true. Yet that rule has no application to the universal emanation of all being from the First Principle, for if the First Principle presupposed anything independent to work upon, that presupposed thing would not be caused by it, contradicting the conclusion that every reality other than the First Being receives its being from it.

Greek metaphysics established participation, the priority of act over potency, the impossibility of an infinite essentially ordered causal regress, and the necessity of an immaterial, immutable First Principle. What neither Plato's Demiurge nor Aristotle's Unmoved Mover reached was the radical contingency of existence itself: the recognition that even eternal matter, were it to exist, would possess no being of its own, and that the deepest composition in every finite thing lies not merely between matter and form, but between what a thing is and the act by which it exists.
